\chapter[TINJAUAN PUSTAKA]{\\ TINJAUAN PUSTAKA}

\section{State Of The Art}
Pembuatan penelitian ini menggunakan referensi dari beberapa jurnal yang berkaitan dan berhubungan. Bentuk referensi ditampilkan melalui tabel ringkasan penelitian sebelumnya atau \textit{State of the Art} yang dapat dilihat pada Tabel \ref{tab:sota} berikut.

\begin{table}[htbp]
    \centering
    \caption{State Of The Art}
    \label{tab:sota}
    \fontsize{9pt}{11pt}\selectfont
    \begin{tabularx}{\textwidth}{|c|>{\raggedright\arraybackslash}p{2.3cm}|>{\raggedright\arraybackslash}X|>{\raggedright\arraybackslash}X|>{\raggedright\arraybackslash}X|>{\raggedright\arraybackslash}X|}
        \hline
        \textbf{No} & \textbf{Penulis, Tahun \& Judul} & \textbf{Metode} & \textbf{Hasil} & \textbf{Persamaan} & \textbf{Perbedaan} \\
        \hline
        1 & Darmanto \& Barus (2024) \cite{Darmanto2024} & Random Forest & Model mampu meningkatkan akurasi prediksi serta membantu perusahaan dalam perencanaan stok. & Sama-sama menggunakan data transaksi penjualan sebagai dasar peramalan. & Menggunakan Random Forest untuk klasifikasi/regresi vs. LSTM untuk runtun waktu. \\
        \hline
        2 & Prasetyo \textit{et al.} (2023) \cite{Prasetyo2023} & LSTM & LSTM mampu menangkap pola tren dan musiman secara lebih baik dibanding regresi linier. & Sama-sama menggunakan metode LSTM untuk peramalan penjualan berbasis data historis. & Hanya fokus pada pengembangan model prediksi vs. integrasi model ke aplikasi POS. \\
        \hline
        3 & Ramadhan \textit{et al.} (2022) \cite{Ramadhan2022} & LSTM dan GRU & LSTM memiliki performa lebih stabil dalam menghadapi fluktuasi permintaan jangka panjang. & Memanfaatkan pendekatan deep learning untuk memodelkan data deret waktu yang kompleks. & Hanya fokus pada analisis model vs. rancang bangun aplikasi POS terintegrasi. \\
        \hline
        4 & Wijaya \& Putri (2022) \cite{Wijaya2022} & - & Sistem POS berbasis web mengurangi kesalahan manual dan mempercepat pelayanan. & Mengembangkan aplikasi POS untuk membantu pengelolaan transaksi dan laporan. & Belum memiliki fitur prediksi penjualan vs. memiliki fitur prediksi LSTM. \\
        \hline
    \end{tabularx}
\end{table}

\section{Tinjauan Teoritis}

\subsection{Point of Sale (POS)}
\textit{Point of Sale} (POS) merupakan sistem informasi yang digunakan untuk memproses dan mencatat transaksi penjualan pada titik terjadinya transaksi antara penjual dan pelanggan \cite{Wijaya2022}. Sistem POS modern telah berkembang dari sekadar mesin kasir elektronik menjadi sistem berbasis komputer atau aplikasi yang terintegrasi dengan basis data dan jaringan internet. Melalui sistem POS, setiap transaksi penjualan akan tercatat secara otomatis dan tersimpan dalam database, sehingga meminimalkan kesalahan pencatatan manual serta meningkatkan kecepatan pelayanan. Selain itu, sistem POS juga umumnya dilengkapi dengan fitur manajemen inventaris yang mampu memperbarui jumlah stok secara otomatis setelah terjadi transaksi, sehingga membantu pelaku usaha dalam mengontrol persediaan barang.

Keunggulan utama sistem POS terletak pada kemampuannya menghasilkan data historis penjualan yang terstruktur dan terdokumentasi dengan baik. Data historis ini menjadi fondasi penting dalam proses analisis bisnis, termasuk analisis tren penjualan, identifikasi produk terlaris, serta peramalan permintaan di masa mendatang. Dengan demikian, sistem POS tidak hanya berperan sebagai alat operasional, tetapi juga sebagai sumber data strategis bagi pengembangan sistem analitik yang lebih canggih.

\subsection{Konsep Peramalan (Forecasting)}
Peramalan atau \textit{forecasting} adalah proses estimasi nilai di masa depan berdasarkan pola data historis dan analisis faktor-faktor yang memengaruhinya \cite{Prasetyo2023}. Dalam konteks bisnis ritel, peramalan penjualan memiliki peran yang sangat penting karena berkaitan langsung dengan perencanaan stok, pengadaan barang, pengelolaan modal, dan strategi pemasaran. Kesalahan dalam peramalan dapat menyebabkan terjadinya kelebihan stok (\textit{overstock}) yang meningkatkan biaya penyimpanan, maupun kekurangan stok (\textit{stockout}) yang berpotensi mengurangi pendapatan dan kepuasan pelanggan.

Metode peramalan dapat dibedakan menjadi metode kualitatif dan kuantitatif. Metode kuantitatif umumnya memanfaatkan data historis numerik dan teknik statistik seperti \textit{Moving Average}, \textit{Exponential Smoothing}, serta \textit{ARIMA}. Namun, seiring berkembangnya teknologi komputasi, pendekatan berbasis \textit{machine learning} dan \textit{deep learning} mulai banyak digunakan karena mampu memodelkan hubungan non-linear yang kompleks serta menangani data dalam jumlah besar \cite{Ramadhan2022}. Pendekatan ini dinilai lebih adaptif terhadap perubahan pola penjualan yang dinamis dibandingkan metode statistik tradisional.

\subsection{Data Deret Waktu (Time Series)}
Data deret waktu (\textit{time series}) merupakan data yang dikumpulkan secara berurutan berdasarkan interval waktu tertentu, seperti harian, mingguan, atau bulanan \cite{Prasetyo2023}. Data penjualan harian termasuk dalam kategori ini karena memiliki dimensi waktu yang berurutan dan saling bergantung. Karakteristik utama data \textit{time series} meliputi tren jangka panjang, pola musiman yang berulang, fluktuasi siklis yang dipengaruhi kondisi ekonomi, serta komponen acak yang tidak dapat diprediksi secara langsung.

Dalam analisis \textit{time series}, penting untuk memahami bahwa nilai pada suatu waktu tertentu dipengaruhi oleh nilai pada waktu sebelumnya. Oleh karena itu, model yang digunakan dalam peramalan harus mampu menangkap ketergantungan temporal tersebut. Pada kasus penjualan harian, misalnya, jumlah penjualan pada hari ini dapat dipengaruhi oleh pola penjualan pada hari-hari sebelumnya, termasuk faktor musiman seperti akhir pekan atau periode promosi tertentu. Pemahaman terhadap karakteristik ini menjadi dasar dalam pemilihan metode LSTM sebagai model prediksi.

\subsection{Machine Learning dalam Prediksi Penjualan}
\textit{Machine Learning} merupakan cabang dari kecerdasan buatan yang memungkinkan sistem belajar dari data untuk mengenali pola dan menghasilkan prediksi tanpa diprogram secara eksplisit \cite{Goodfellow2016}. Dalam prediksi penjualan, pendekatan \textit{supervised learning} digunakan karena tersedia data historis yang berisi pasangan input dan output yang telah diketahui nilainya. Model akan dilatih menggunakan data tersebut untuk mempelajari pola hubungan antara waktu dan nilai penjualan.

Metode \textit{machine learning} konvensional seperti regresi linier atau \textit{Random Forest} dapat digunakan untuk memprediksi penjualan, namun metode tersebut memiliki keterbatasan dalam menangani data sekuensial yang memiliki ketergantungan jangka panjang. Oleh karena itu, pendekatan berbasis jaringan saraf tiruan, khususnya \textit{Recurrent Neural Network} (RNN), dikembangkan untuk mengatasi keterbatasan tersebut dengan memanfaatkan struktur jaringan yang mampu memproses data berurutan.

\subsection{Recurrent Neural Network (RNN)}
\textit{Recurrent Neural Network} (RNN) merupakan arsitektur jaringan saraf tiruan yang dirancang untuk memproses data sekuensial dengan mempertahankan informasi dari langkah waktu sebelumnya \cite{Goodfellow2016}. RNN memiliki mekanisme umpan balik (\textit{feedback loop}) yang memungkinkan output pada waktu sebelumnya menjadi input tambahan pada waktu berikutnya. Struktur ini memungkinkan model memahami konteks temporal dalam data deret waktu.

Meskipun RNN mampu memodelkan data sekuensial, arsitektur ini memiliki kelemahan berupa permasalahan \textit{vanishing gradient}, yaitu kondisi di mana gradien menjadi sangat kecil sehingga jaringan kesulitan mempelajari hubungan jangka panjang. Permasalahan ini menyebabkan performa RNN menurun ketika digunakan untuk memodelkan data dengan dependensi waktu yang panjang, seperti data penjualan harian dalam periode yang cukup lama.

\subsection{Long Short-Term Memory (LSTM)}
\textit{Long Short-Term Memory} (LSTM) merupakan salah satu varian dari \textit{Recurrent Neural Network} (RNN) yang dirancang untuk mengatasi permasalahan \textit{vanishing gradient} pada pemodelan data sekuensial jangka panjang \cite{Hochreiter1997}. LSTM diperkenalkan oleh Sepp Hochreiter dan Jürgen Schmidhuber pada tahun 1997 sebagai arsitektur jaringan saraf yang mampu mempertahankan informasi dalam rentang waktu yang lebih panjang dibandingkan RNN konvensional. Keunggulan utama LSTM terletak pada keberadaan mekanisme \textit{gating} (gerbang) yang mengatur aliran informasi masuk, disimpan, dan dikeluarkan dari sel memori (\textit{cell state}). Struktur ini memungkinkan model untuk memilih informasi mana yang relevan untuk dipertahankan atau dilupakan, sehingga sangat efektif dalam berbagai tugas seperti peramalan deret waktu, pengenalan suara, dan pemrosesan bahasa alami \cite{Goodfellow2016}.

Secara arsitektural, satu unit LSTM terdiri atas tiga gerbang utama, yaitu \textit{forget gate}, \textit{input gate}, dan \textit{output gate}, serta satu komponen \textit{cell state} yang berfungsi sebagai memori jangka panjang. Pada setiap langkah waktu ke-$t$, LSTM menerima input $x_t$, status tersembunyi sebelumnya $h_{t-1}$, dan status sel sebelumnya $C_{t-1}$. Proses komputasi dilakukan secara bertahap melalui fungsi aktivasi sigmoid dan tanh untuk mengontrol pembaruan memori. Pendekatan ini memungkinkan LSTM menjaga stabilitas gradien selama proses pelatihan menggunakan algoritma \textit{backpropagation through time} (BPTT). Dengan demikian, LSTM mampu menangkap dependensi temporal yang kompleks dalam data sekuensial.

Secara matematis, mekanisme kerja LSTM dapat dirumuskan sebagai berikut:

\begin{enumerate}
    \item \textbf{Forget Gate} \\
    Gerbang ini menentukan informasi mana dari status sel sebelumnya yang akan dipertahankan atau dihapus.
    \begin{equation}
        f_t = \sigma(W_f \cdot [h_{t-1}, x_t] + b_f)
    \end{equation}
    Di mana $f_t$ adalah nilai \textit{forget gate} pada waktu ke-$t$, $W_f$ adalah matriks bobot \textit{forget gate}, $b_f$ adalah bias, dan $\sigma$ adalah fungsi aktivasi sigmoid.

    \item \textbf{Input Gate} \\
    Gerbang ini menentukan informasi baru yang akan disimpan dalam \textit{cell state}.
    \begin{equation}
        i_t = \sigma(W_i \cdot [h_{t-1}, x_t] + b_i)
    \end{equation}
    Kandidat nilai memori baru dihitung dengan:
    \begin{equation}
        \tilde{C}_t = \tanh(W_c \cdot [h_{t-1}, x_t] + b_c)
    \end{equation}
    Di mana $i_t$ adalah nilai \textit{input gate}, dan $\tilde{C}_t$ adalah kandidat pembaruan memori.

    \item \textbf{Pembaruan Cell State} \\
    Pembaruan \textit{cell state} lama menjadi \textit{cell state} baru dihitung dengan mengombinasikan \textit{forget gate} dan \textit{input gate}:
    \begin{equation}
        C_t = f_t * C_{t-1} + i_t * \tilde{C}_t
    \end{equation}
    Di mana $C_t$ adalah \textit{cell state} baru, dan $C_{t-1}$ adalah \textit{cell state} sebelumnya.

    \item \textbf{Output Gate} \\
    Gerbang ini menentukan status tersembunyi (\textit{hidden state}) berikutnya yang akan ditransmisikan:
    \begin{equation}
        o_t = \sigma(W_o \cdot [h_{t-1}, x_t] + b_o)
    \end{equation}
    Status tersembunyi $h_t$ diperoleh melalui:
    \begin{equation}
        h_t = o_t * \tanh(C_t)
    \end{equation}
    Di mana $o_t$ adalah nilai \textit{output gate}, dan $h_t$ adalah \textit{hidden state} pada waktu ke-$t$.
\end{enumerate}

Melalui kombinasi mekanisme gerbang tersebut, LSTM mampu mempertahankan informasi penting dalam jangka panjang sekaligus mengurangi pengaruh informasi yang tidak relevan. Secara empiris, arsitektur ini menunjukkan performa yang unggul dalam berbagai studi komparatif terhadap RNN standar, khususnya pada data dengan pola temporal kompleks dan non-linear \cite{Goodfellow2016}.

\subsection{Evaluasi Kinerja Model}
Evaluasi kinerja model prediksi bertujuan untuk mengukur tingkat kesalahan antara nilai aktual dan nilai hasil prediksi yang dihasilkan oleh model \cite{Prasetyo2023}. Dalam penelitian ini, evaluasi dilakukan menggunakan metrik kuantitatif yang umum digunakan pada data \textit{time series}, yaitu \textit{Mean Absolute Error} (MAE), \textit{Root Mean Squared Error} (RMSE), dan \textit{Mean Absolute Percentage Error} (MAPE).

\begin{enumerate}
    \item \textbf{Mean Absolute Error (MAE)} \\
    MAE mengukur rata-rata selisih absolut antara nilai aktual dan prediksi.
    \begin{equation}
        \text{MAE} = \frac{1}{n} \sum_{i=1}^{n} |y_i - \hat{y}_i|
    \end{equation}
    
    \item \textbf{Root Mean Squared Error (RMSE)} \\
    RMSE memberikan penalti lebih besar terhadap kesalahan yang signifikan karena menggunakan kuadrat dari selisih tersebut.
    \begin{equation}
        \text{RMSE} = \sqrt{\frac{1}{n} \sum_{i=1}^{n} (y_i - \hat{y}_i)^2}
    \end{equation}

    \item \textbf{Mean Absolute Percentage Error (MAPE)} \\
    MAPE digunakan untuk mengukur kesalahan dalam bentuk persentase sehingga lebih mudah diinterpretasikan dalam konteks bisnis.
    \begin{equation}
        \text{MAPE} = \frac{100\%}{n} \sum_{i=1}^{n} \left| \frac{y_i - \hat{y}_i}{y_i} \right|
    \end{equation}
\end{enumerate}

Proses evaluasi ini penting untuk memastikan bahwa model LSTM yang diintegrasikan dalam aplikasi POS memiliki tingkat akurasi yang memadai dan dapat diandalkan dalam mendukung pengambilan keputusan operasional.